% BOOK-FILE-ID: f-494e45749d0a
\subsection{目指す形}\label{節：目指す形}
  これから先で扱う数々のパターンと,それから作られるさらに複雑な漸化式,ひいては数学の問題を解くという過程の中で通底する,重要な発想や知識をここでは紹介しよう.
  それらは大きく分けて2つである.
  \begin{enumerate}
    \item 実験をする\\
          $\rightarrow$ シンプルな場合で実験\\
          $\rightarrow$ 答えを予測
    \item 既知の形をつくる\\
          $\rightarrow$ 日々の演習\\
          $\rightarrow$ 勘を鍛える
  \end{enumerate}
  
  こと漸化式や整数問題においては,特に実験が重要である.
  小さな値での振る舞いを見ることで,一般項を予測できたり,大きな値での振る舞いを見ることで,数列全体の流れを掴むことができる.
  
  既知の形に帰着させられれば,あとは計算を行うだけで済む.
  だから重要なのはいかにして知っている形を作れるか,言い換えれば"きれい"な形が作れるかである.
  ここで言う"きれい"というのは,演習を重ねるにつれて理解される非言語の勘ではあるが,あえて言語化すれば
  \begin{enumerate}
    \item 等比型になっている
    \item 定数型$b_{n+1}=b_n$になっている
    \item 左辺が$n+1$に関する式,右辺が$n$に関する同じ形の式になっている
  \end{enumerate}
  などである.定数型は各項が同じ値を取る形であり,\sectionref*{節：定数型}では補助数列を置いてこの形を作る.
  
  既知のものを増やしていくためには,きちんと日々演習を積むことが重要である.
  演習によって「この数字を見たことがあるな」とか「この関数見たことがあるな」とかそういったことを思えるようになる.
  その勘が解法の素早い選択のためには必要になる.

  うまく式変形をして既知の形に帰着させて一般項を得ることだけが手立てではない.
  数学的帰納法を使って一般項を示すのも立派な方法である.\footnote{数学的帰納法については\sectionref{節：数学的帰納法}を参照.}
  一般項を予測して数学的帰納法で一般項を示す手法は,「きっと一般項はこれだ！」という気づきを得る必要がある.
  その点で一定のセンスや経験が必要だろう.
  だが,できないとやらないは大きな違いがある.
  できるようになっておこう.

  \begin{bookpoint}
    \textbf{POINT}\par
    数項を計算して変化を観察し,既知の形を作る置き換えを探す.
    一般項を予想したら,初期条件と漸化式を使ってすべての項で成り立つことを示す.
  \end{bookpoint}
